\documentclass[11pt]{article}

\usepackage[a4paper,margin=1in]{geometry}
\usepackage{amsmath,amssymb}
\usepackage{hyperref}

\hypersetup{
    colorlinks=true,
    linkcolor=blue,
    citecolor=blue,
    urlcolor=blue
}

\title{Proof That the First Variation of Entropy Is $\log\rho+1$}
\author{}
\date{}

\begin{document}

\maketitle

\section*{Statement}

Let
\[
\mathcal H(\rho)
=
\int_{\mathbb R^d}\rho(x)\log\rho(x)\,dx
\]
be the Boltzmann entropy functional, using the convention common in
Wasserstein gradient flows.
Assume $\rho$ is positive and smooth enough for the following computation, and
let $\eta$ be a smooth perturbation satisfying
\[
\int_{\mathbb R^d}\eta(x)\,dx=0.
\]
Then
\[
\frac{\delta \mathcal H}{\delta \rho}
=
\log\rho+1.
\]

\section*{Proof}

Consider the perturbed density
\[
\rho_\varepsilon=\rho+\varepsilon\eta.
\]
For sufficiently small $\varepsilon$, this remains positive if $\rho$ is
strictly positive and $\eta$ is bounded. We compute
\[
\left.
\frac{d}{d\varepsilon}
\mathcal H(\rho_\varepsilon)
\right|_{\varepsilon=0}.
\]
By definition,
\[
\mathcal H(\rho_\varepsilon)
=
\int_{\mathbb R^d}
(\rho(x)+\varepsilon\eta(x))
\log(\rho(x)+\varepsilon\eta(x))\,dx.
\]
Differentiate the scalar function
\[
h(r)=r\log r.
\]
Its derivative is
\[
h'(r)=\log r+1.
\]
Therefore, pointwise in $x$,
\[
\left.
\frac{d}{d\varepsilon}
\Big[
(\rho(x)+\varepsilon\eta(x))
\log(\rho(x)+\varepsilon\eta(x))
\Big]
\right|_{\varepsilon=0}
=
(\log\rho(x)+1)\eta(x).
\]
Integrating gives
\[
\left.
\frac{d}{d\varepsilon}
\mathcal H(\rho+\varepsilon\eta)
\right|_{\varepsilon=0}
=
\int_{\mathbb R^d}
(\log\rho(x)+1)\eta(x)\,dx.
\]
By the definition of first variation, this means
\[
\frac{\delta \mathcal H}{\delta \rho}(\rho)(x)
=
\log\rho(x)+1.
\]

\section*{Remark on the Constant}

When working on probability densities, admissible perturbations satisfy
\[
\int_{\mathbb R^d}\eta(x)\,dx=0.
\]
Hence adding any constant $C$ to the first variation does not change the
directional derivative, because
\[
\int_{\mathbb R^d}C\,\eta(x)\,dx=0.
\]
Thus, on the probability simplex, the first variation is defined only up to an
additive constant. One often writes simply
\[
\frac{\delta \mathcal H}{\delta\rho}=\log\rho
\]
when the constant is irrelevant. In Wasserstein gradient flows, the constant
also disappears after taking a spatial gradient:
\[
\nabla(\log\rho+1)=\nabla\log\rho.
\]

\end{document}
